%
% File: concl.tex
%
\chapter{Conclusion}
\setlength{\parskip}{1em}

This work presented the design, implementation, and operational preparation of an intelligent meal planning system for Kazakhstani households. The work combined five sibling codebases, a multi-source ETL pipeline producing a published catalog of approximately ninety thousand recipes, a three-tier hybrid planner that joins a deterministic candidate-pool builder, an OR-Tools CP-SAT combinatorial selection, and a LightGBM listwise scoring layer with a fallback chain, a React Native mobile client structured around twenty-seven screens, a React + Vite web client, and a Terraform-managed AWS deployment plan with eight GitHub OIDC workflows and five operational runbooks. Five architecture decision records justify the lasting choices: monorepo, Fargate, hybrid CP-SAT plus ML, Expo Managed, and Postgres. Evaluation reports concrete numbers from the repository state and explicitly distinguishes deployment-ready capabilities from claimed production performance.

The central argument is that disciplined software engineering, classical operations research, and modern learning-to-rank can be combined end-to-end within a single decision-support application. Constraint-based household planning benefits from a global combinatorial model rather than from greedy local choice; learning-to-rank is an appropriate but supporting extension; production engineering matters as much as algorithmic elegance, because a system that cannot be deployed, monitored, and recovered is not a system. The accompanying repository codifies these commitments through versioned data, reproducible bootstrap commands, infrastructure-as-code, and runbooks, so that the work continues to be inspectable and extensible.

Returning to the five research questions stated in the Introduction: a multi-source ETL pipeline with embedding-based clustering, MinHash deduplication, USDA-derived nutrition synchronisation, and an eight-component quality score answers RQ1 by producing a reproducibly catalogued and quality-gated corpus of more than ninety thousand recipes; a three-tier hybrid planner combining deterministic prefiltering, OR-Tools CP-SAT assignment, and a learning-to-rank scoring layer with graceful fallback answers RQ2 by reducing macro-target gaps from thousands of kilocalories and hundreds of grams of protein down to noise on the reported scenario; a five-codebase architecture coordinated by a versioned API and per-directory CI/CD answers RQ3 by isolating the algorithmic core from delivery concerns while keeping a single OpenAPI contract for both mobile and web clients; a Terraform-managed AWS deployment with GitHub OIDC, observability, security middleware, and five runbooks answers RQ4 by satisfying recoverability, security, and store-submission requirements; and an evaluation strategy that distinguishes deployment-ready capabilities from production-measured performance answers RQ5 by reporting only honest, reproducible numbers.

Future work falls naturally into four staged tracks. In the data track, increased canonical ingredient coverage, reduced manual nutrition curation through improved alias learning, and an Astana-specific Arbuz snapshot extend the existing baseline. In the planner track, pantry-aware scoring, leftover-aware planning, per-member portion scaling, hard calorie-target bands, and Pareto-frontier plan generation refine the optimisation surface. In the machine learning track, online evaluation against the deterministic baseline through PostHog feature flags, feature-importance drift detection, and decile-bucket calibration unlock measured performance gains once production feedback accumulates. In the production track, the work labelled here as ``ready'' becomes ``deployed'' upon the team activating an AWS account and submitting to the App Store and Google Play. The accompanying runbooks, ADRs, and infrastructure-as-code make that transition a business decision rather than an engineering blocker.

\clearpage
